% =============================================================================
%  QuPath_Rust preprint -- arXiv (eess.IV, cross-list cs.CV)
%
%  Build:  latexmk -pdf main.tex
%     or:  pdflatex main && bibtex main && pdflatex main && pdflatex main
%
%  STATUS: skeleton. Every section is an outline with the evidence it must
%  carry named in a \todo{}. Nothing here is written prose yet.
%
%  The argument and most of the numbers already exist in the repository, in
%  DIVERGENCES.md (D-009) and docs/MEMO.md section 2. Drafting this is largely
%  restructuring those, not new analysis. Copy numbers from there rather than
%  re-deriving them, and cite the reproduction command for each.
% =============================================================================
\documentclass[11pt,a4paper]{article}

\usepackage[utf8]{inputenc}
\usepackage[T1]{fontenc}
\usepackage{amsmath,amssymb}
\usepackage{graphicx}
\usepackage{booktabs}
\usepackage{xcolor}
\usepackage[margin=1in]{geometry}
\usepackage[hidelinks]{hyperref}

% Draft notes. Set \draftnotesfalse for a clean read-through.
\newif\ifdraftnotes
\draftnotestrue
\ifdraftnotes
  \newcommand{\todo}[1]{{\color{red}\small[\textbf{TODO:} #1]}}
\else
  \newcommand{\todo}[1]{}
\fi

\title{Verifying a port of a digital-pathology pipeline,\\
       stage by stage}

% TODO: real name, affiliation and ORCID. The repository currently carries
% only the handle "Dante-Berth", which is not submittable.
\author{\todo{author, affiliation, ORCID}}
\date{\today}

\begin{document}
\maketitle

\begin{abstract}
\todo{Write last. The frame: a port of QuPath's watershed cell detection to
Rust, and --- the actual contribution --- a stage-boundary fixture method for
verifying such a port against the reference implementation. Say plainly in the
abstract that a density-dependent count difference of 1.065x remains open and
is localised to three of ten stages. Claiming parity here would be refuted by
the repository's own DIVERGENCES.md, which a reviewer will read.}
\end{abstract}

\section{Introduction}
\label{sec:intro}

\todo{QuPath's role in digital pathology \cite{bankhead2017qupath}, and the
bulk-analysis bottleneck: a per-object string-keyed measurement map is the
wrong layout for batch work over many slides.}

\todo{The real problem: a port must \emph{prove} it computes the same thing
before anyone can use its output in research. State the thesis here --- that
agreeing final object counts are the weakest available evidence, because two
runs can agree on how many objects they found and disagree about where every
one of them is.}

\todo{Contributions, as a list: (i) the stage-boundary fixture method;
(ii) mutation testing applied to the fixtures themselves; (iii) three bugs
found that produced entirely plausible images; (iv) a localisation result for
the remaining discrepancy.}

\section{The verification method}
\label{sec:method}

\todo{The core idea, and the one figure worth drawing from scratch: record the
reference implementation's intermediate state at each stage boundary, feed the
port the \emph{same} input at that boundary, and compare pixel for pixel. This
isolates each stage --- a stage that faithfully processes a different input
produces different output, so end-to-end comparison cannot tell you which
stage is wrong.}

\todo{Recording: parity/java/*.java run inside a live QuPath build. Say which
recorders need a slide (parity/README.md documents this).}

\todo{Mutation testing \cite{jia2021mutation} as the check on the checks --- the
part most ports skip. A parity test that cannot fail proves nothing, so every
one here was shown to fail under a deliberate perturbation before being
trusted. Give the concrete examples: dropping the hole fill (246 pixels),
narrowing the dilation to five neighbours (97 pixels), moving the maximum
finder's tolerance from 0.5 to 0.6 (55 pixels).}

\todo{Tolerances, and why they are stated per fixture rather than globally.}

\section{The port}
\label{sec:port}

\todo{The columnar object layout: a struct of arrays, one contiguous column
per feature, one row per object. Be precise, as the README already is, that
this --- not the choice of language --- is the main source of the speedup. Do
not let the paper lose that distinction.}

\todo{The detection chain, following QuPath's WatershedCellDetection: colour
deconvolution \cite{ruifrok2001quantification}, background estimation by
morphological reconstruction \cite{vincent1993morphological}, Laplacian of
Gaussian, regional maxima, labelling, watershed flood
\cite{vincent1991watersheds}, region tracing, the mean-intensity filter, the
merge, the shape split, refineBoundary, the area filter.}

\todo{Slide reading: a pure-Rust Aperio SVS/tiled-TIFF reader, checked against
OpenSlide \cite{goode2013openslide}.}

\todo{What is deliberately \emph{not} claimed: the viewer claims no parity and
no speedup, and smooth\_boundary / expand\_labels / cell\_areas are ported and
pinned but uncalled. Saying so costs nothing and is what makes the rest
credible.}

\section{Results: what the method found}
\label{sec:results}

\subsection{Three bugs behind plausible images}
\todo{The strongest argument for the method, because visual inspection passed
all three: (i) a morphological opening in place of a top hat --- wrong
operation, wrong radius (2px vs 16), wrong direction; (ii) the DAB stain vector
used under the name EOSIN, a 52\% error on the leading coefficient; (iii)
Aperio JPEG tiles carrying no Adobe APP14 marker, so a conforming decoder
picks the wrong colour transform.}

\subsection{Performance}
\todo{Reuse docs/benchmarks/compare-threads.svg. Carry its caveat: the
4000x4000 region is four tiles and saturates at four threads; the 9000x9000
region is twenty-five tiles and keeps improving to twenty. A stage share
quoted without its tile and thread count is not a number anyone can use.}

\subsection{Object-level agreement}
\todo{The Dice table over three single-tile regions. The point that pooled Dice
involves no matching decision, so it is read first; and that both sides are the
pixels each pipeline labelled, not a rasterisation of our polygon, which would
have measured our rasteriser instead.}

\begin{table}[t]
  \centering
  \caption{Nucleus-level agreement over three single-tile regions of CMU-1.
           \todo{copy from DIVERGENCES.md D-009; add the caption's reading}}
  \label{tab:dice}
  \begin{tabular}{lrrrr}
    \toprule
    region & ratio & pooled Dice & matched at Dice $\geq 0.5$ & unmatched \\
    \midrule
    \todo{fill from D-009} & & & & \\
    \bottomrule
  \end{tabular}
\end{table}

\section{The open gap}
\label{sec:gap}

\todo{State it plainly and early: 1.065x pooled over ten regions of CMU-1,
156{,}979 objects against 147{,}351 cells, 1.012--1.129 by region, always
over-counting. This is the paper's honest centre, not a limitation buried at
the end.}

\todo{What it is not. Seven stages ruled out against QuPath on real tissue,
pixel for pixel: the merge, the shape split, the colour deconvolution, the
JPEG decode beneath it, the labelling, the flood, and the background estimate
--- the last at the 2048x2048 the pipeline actually runs it at, after the
160x160 fixture was shown to be too small to see the effect.}

\todo{The refuted hypothesis, kept because it is a result: the merge/split pair
was the named suspect, and MergeAndSplitTile.java exonerated it --- the merged
mask, MaximumFinder's segmentation and the mask toWatershed leaves are
identical pixel for pixel on all eight crops, in both density regimes. The
stages where a gap is \emph{visible} are not the stages where it is
\emph{made}.}

\todo{The stage ladder, which is the strongest single table: it splits the gap
in two. A flat, density-\emph{independent} $\sim$4.7\% excess already present
in the regional maxima, and a separate density-\emph{dependent} jump appearing
at the mean-intensity filter (sparse 1.149 against dense 1.049 and 1.040).
Conflating these was what the too-small crops forced.}

\begin{table}[t]
  \centering
  \caption{Object counts at four stage boundaries, QuPath $\rightarrow$ this
           port. \todo{copy from D-009}}
  \label{tab:ladder}
  \begin{tabular}{lrrr}
    \toprule
    rung & sparse & dense-2 & dense \\
    \midrule
    \todo{fill from D-009} & & & \\
    \bottomrule
  \end{tabular}
\end{table}

\todo{The clumping measurement, read off the image rather than inferred:
maxima per thresholded-LoG blob is 1.84 on the sparse region against 1.46 and
1.42 on the dense ones.}

\todo{Frame the section's conclusion as a localisation result: three stages
remain --- the blur, the Laplacian and the regional maxima --- each with a
synthetic fixture and none with one on tissue at the size it runs at. That
pattern (agreement on constructed input, untested on real input at production
size) is the one this project has repeatedly found worth distrusting.}

\section{Discussion}
\label{sec:discussion}

\todo{What transfers to other scientific-software ports: stage-boundary
fixtures, mutation-checked; recording negative results; and the specific
lesson that a fixture must run at the size the pipeline gives the stage, since
an operation switching strategy on a global condition can agree at one size
and diverge at another.}

\todo{Limits, stated without hedging: one slide (CMU-1), one detection
algorithm, one reference version of QuPath. State CMU-1's provenance and
licence --- it is from the OpenSlide freely-distributable test set --- since
every headline number depends on it.}

\todo{Future work: closing D-009 by pinning the three remaining stages on
tissue at production size.}

\section*{Data and code availability}
\todo{Repository URL, the archived release DOI (Zenodo), the licence
(GPL-3.0-or-later, matching QuPath as a derivative work), and the slide's
source. Name the reproduction command for each headline number, as the README
already does.}

\section*{Acknowledgements}
\todo{The QuPath authors. Notify them before submitting --- this is a
derivative work reimplementing their algorithms, and a reviewer may well be
one of them.}

\bibliographystyle{plain}
\bibliography{../paper}

\end{document}
